\subsection{window-\texorpdfstring{$6$}{6} 的极小 Lie 完成、自旋超立方体与 dyadic 混合签名}\label{subsec:conclusion-window6-spin-minuscule-hypercube-signature}
在 Fibonacci--Lie 共振已经把 \(21=\dim \mathfrak{so}(7)\) 外部化、分辨率塔的 Haar 归一化与 \(2\)-进秩二刚性已经把 \(T_2(E_{\min})\cong \ZZ_2^2\) 固定、而 window-\(6\) 的 root--Cartan 完成与 boundary parity 识别也已闭合之后，还可进一步把“极小 Lie 型如何被最小纤维层选出”与“连续圆维如何同 dyadic 地址账本耦合”压缩为同一条三层链。这里所谓 Lie 完成始终指与 \(X_6=X_6^{\mathrm{cyc}}\sqcup X_6^{\mathrm{bdry}}\) 的 root--Cartan 字典相容的紧连续包络，而非把有限集合 \(X_6\) 本身视为某个连通 Lie 群的连续轨道空间。

\begin{theorem}[window-\texorpdfstring{$6$}{6} 的极小紧 Lie 完成与圆维刚性]\label{thm:conclusion-window6-minimal-compact-lie-completion-circle-rigidity}
\leanverified{paper\_conclusion\_window6\_minimal\_compact\_lie\_completion\_circle\_rigidity}
令
\[
\mathcal V_6:=\RR\langle X_6\rangle.
\]
则与 \(X_6=X_6^{\mathrm{cyc}}\sqcup X_6^{\mathrm{bdry}}\) 的 root--Cartan 字典相容的极小紧单 Lie 代数完成，在忘记 boundary 排序后的同构意义下，恰有两类：
\[
\mathfrak k_6^{(B)}\cong \mathfrak{so}(7),
\qquad
\mathfrak k_6^{(C)}\cong \mathfrak{sp}(3).
\]
更精确地，\(X_6^{\mathrm{cyc}}\) 必对应于 \(18\) 个根方向，\(X_6^{\mathrm{bdry}}\) 必对应于 \(3\) 维 Cartan 子代数；因此任一与该字典兼容的极小紧连通 Lie 群完成 \(K_6\) 都满足
\[
\dim K_6=21,
\qquad
\operatorname{rank}(K_6)=3.
\]
若以其极大环面 \(T_6\) 的圆维定义 window-\(6\) 的连通可见圆维，则
\[
\cdim(T_6)=3.
\]
\end{theorem}
\begin{proof}
定理 \ref{thm:conclusion-window6-abstract-chevalley-completion} 已给出两个 \(21\) 维实 Lie 代数完成，其复化分别为 \(B_3\) 型与 \(C_3\) 型。另一方面，定理 \ref{thm:window6-21-adjoint-weight-multiset} 已把 \(X_6\) 规范识别为 \(18\) 条根加 \(3\) 重零权；因此任一与该 root--Cartan 字典相容的极小不可约完成都必须来自一个秩 \(3\)、根数 \(18\) 的约化不可约根系。秩 \(3\) 情形中只有 \(B_3\) 与 \(C_3\) 满足这一计数，故 Lie 型只能是这两类。

取其紧实型即得 \(\mathfrak{so}(7)\) 与 \(\mathfrak{sp}(3)\)。于是相应紧连通 Lie 群完成满足
\[
\dim K_6=18+3=21,
\qquad
\operatorname{rank}(K_6)=3.
\]
最后由定理 \ref{thm:conclusion-window6-boundary-torus-2torsion}，其极大环面 \(T_6\cong (S^1)^3\)，故
\[
\cdim(T_6)=3.
\]
证毕。
\end{proof}

\begin{theorem}[八个最小纤维唯一锁定 \texorpdfstring{$B_3$}{B3}，从而锁定 \texorpdfstring{$\mathrm{Spin}(7)$}{Spin(7)}]\label{thm:conclusion-window6-minuscule-eightfold-b3-spin7}
\leanverified{paper\_conclusion\_window6\_minuscule\_eightfold\_b3\_spin7}
沿用定理 \ref{thm:xi-window6-center-exact-splitting-bdry-cyclic} 的记号，记
\[
Y_2:=X_6^{(2)}=\{w\in X_6:\ d_6^{\mathrm{bin}}(w)=2\}.
\]
则 \(|Y_2|=8\)。若进一步要求这一最小纤维层在定理 \ref{thm:conclusion-window6-minimal-compact-lie-completion-circle-rigidity} 的同一极小 Lie 完成的某个辅助非零 minuscule 表示中可被实现为单一 Weyl 权轨道，则 \(B_3/C_3\) 的二义性被唯一消去：只剩 \(B_3\) 型。等价地，其单连通紧完成唯一为
\[
K_{6,\mathrm{sc}}\cong \mathrm{Spin}(7).
\]
并且存在一个 Weyl-等变识别，把 \(Y_2\) 与 \(\mathrm{Spin}(7)\) 自旋表示的八个权
\[
\left\{
\left(\pm\frac12,\pm\frac12,\pm\frac12\right)
\right\}
\]
对应起来。
\end{theorem}
\begin{proof}
由定理 \ref{thm:xi-window6-center-exact-splitting-bdry-cyclic}，window-\(6\) 的最小纤维层正好由 \(8\) 个稳定类型组成。另一方面，定理 \ref{thm:conclusion-window6-minimal-compact-lie-completion-circle-rigidity} 已表明极小 Lie 型只可能是 \(B_3\) 或 \(C_3\)。

现在加入“\(Y_2\) 应当在同一完成的某个辅助表示中形成单一非零 minuscule Weyl 轨道”的约束。对 \(B_3\)，唯一的非零 minuscule 最高权是自旋权，其权集恰有 \(8\) 个元素；对 \(C_3\)，唯一的非零 minuscule 最高权给出的标准表示只有 \(6\) 个权。故 \(|Y_2|=8\) 这一计数只可能由 \(B_3\) 满足，而 \(C_3\) 被排除。

因此极小 Lie 完成唯一落在 \(B_3\) 侧；其单连通紧群即为 \(\mathrm{Spin}(7)\)。而 \(B_3\) 自旋表示的权集正是
\[
\left\{
\left(\pm\frac12,\pm\frac12,\pm\frac12\right)
\right\},
\]
且 minuscule 权集本身就是一条单一 Weyl 轨道，故 \(Y_2\) 与该八点权集之间存在 Weyl-等变识别。证毕。
\end{proof}

\begin{corollary}[异常电荷格的 toral 商与圆维的模二恢复]\label{cor:conclusion-window6-anomaly-toral-quotient-mod2-recovery}
\leanverified{paper\_conclusion\_window6\_anomaly\_toral\_quotient\_mod2\_recovery}
记
\[
A_6:=\bigl(\mathcal G_6^{\mathrm{bin}}\bigr)^{\mathrm{ab}}
\cong
(\ZZ/2\ZZ)^{21}.
\]
则存在规范分裂短正合列
\[
0\longrightarrow A_6^{\mathrm{root}}
\longrightarrow A_6
\longrightarrow T_6[2]
\longrightarrow 0,
\]
其中
\[
A_6^{\mathrm{root}}\cong (\ZZ/2\ZZ)^{18},
\qquad
T_6[2]\cong (\ZZ/2\ZZ)^3.
\]
因此，window-\(6\) 的 \(21\) 维异常奇偶荷并不处于同一层级：其中 \(18\) 维是根方向通道的模二阴影，\(3\) 维是环面相位的模二阴影；并且极小紧完成的连通圆维可由
\[
\cdim(T_6)=\dim_{\FF_2}T_6[2]=3
\]
精确恢复。
\end{corollary}
\begin{proof}
定理 \ref{thm:window6-abelianized-parity-charge-root-cartan-splitting} 已给出规范直和分解
\[
A_6
\cong
\FF_2[\Phi_{\mathrm{short}}]
\oplus
\FF_2[\Phi_{\mathrm{long}}]
\oplus
\mathfrak h_{\FF_2}.
\]
定义
\[
A_6^{\mathrm{root}}
:=
\FF_2[\Phi_{\mathrm{short}}]\oplus \FF_2[\Phi_{\mathrm{long}}],
\]
便得到一个规范分裂短正合列
\[
0\longrightarrow A_6^{\mathrm{root}}
\longrightarrow A_6
\longrightarrow \mathfrak h_{\FF_2}
\longrightarrow 0.
\]
再由定理 \ref{thm:conclusion-window6-boundary-torus-2torsion}，\(\mathfrak h_{\FF_2}\) 正可识别为 \(T_6[2]\)。维数公式
\[
\dim_{\FF_2}A_6^{\mathrm{root}}=6+12=18,
\qquad
\dim_{\FF_2}T_6[2]=3
\]
遂直接得到。最后，由定理 \ref{thm:conclusion-window6-minimal-compact-lie-completion-circle-rigidity} 中 \(T_6\cong (S^1)^3\)，有
\[
\cdim(T_6)=3=\dim_{\FF_2}T_6[2].
\]
证毕。
\end{proof}

\begin{theorem}[最小纤维层的自旋权超立方体与 boundary 符号作用]\label{thm:conclusion-window6-minimal-fiber-spin-hypercube}
\leanverified{paper\_conclusion\_window6\_minimal\_fiber\_spin\_hypercube}
在定理 \ref{thm:conclusion-window6-minuscule-eightfold-b3-spin7} 的 \(\mathrm{Spin}(7)\) 完成下，\(Y_2\) 可 Weyl-等变识别为自旋表示的权集
\[
\left\{
\left(\pm\frac12,\pm\frac12,\pm\frac12\right)
\right\}.
\]
若两点仅在一个坐标上变号，则定义其相邻。则 \(Y_2\) 的邻接图同构于三维超立方体 \(Q_3\)。

进一步，boundary 三个中心 involution 通过
\[
T_6[2]\cong (\ZZ/2\ZZ)^3
\]
在自旋模的各个权空间上按特征标对角作用；因而在适当权基下，它们可同时实现为三组彼此交换的符号矩阵。于是最小纤维层不再只是抽象八点，而是一个由自旋权与 boundary 二挠共同组织的超立方体层。
\end{theorem}
\begin{proof}
由定理 \ref{thm:conclusion-window6-minuscule-eightfold-b3-spin7}，\(Y_2\) 已可与 \(\mathrm{Spin}(7)\) 自旋表示的八个权作 Weyl-等变识别。该权集恰为
\[
\left\{
\left(\pm\frac12,\pm\frac12,\pm\frac12\right)
\right\},
\]
而“只改动一个坐标符号”的关系正是三维立方体顶点集的自然邻接关系，故所得邻接图即为 \(Q_3\)。

再令 \(S_{\mathrm{spin}}\) 表示相应的八维自旋模，则其按权分解为
\[
S_{\mathrm{spin}}
=
\bigoplus_{\lambda\in Y_2} S_\lambda,
\qquad
\dim_\CC S_\lambda=1.
\]
由定理 \ref{thm:conclusion-window6-boundary-torus-2torsion}，boundary 三个中心 involution 与 \(T_6[2]\) 的三个生成元对应。由于 \(T_6[2]\subset T_6\) 位于极大环面中，每个 \(t\in T_6[2]\) 在权空间 \(S_\lambda\) 上都以标量 \(\lambda(t)\in\{\pm1\}\) 作用。选取与上述权分解相容的基后，这三个 involution 便同时成为三组彼此交换的对角符号矩阵。证毕。
\end{proof}

\begin{proposition}[圆维--dyadic 混合签名的极小性]\label{prop:conclusion-window6-dyadic-mixed-signature-minimality}
令
\[
P_2:=T_2(E_{\min})\cong \ZZ_2^2,
\qquad
\mathcal E_6^{\mathrm{mix}}:=T_6\times P_2.
\]
对形如 \(T\times P\) 的对象，其中 \(T\) 为紧环面、\(P\) 为 pro-\(2\) 交换群，定义其 dyadic 混合签名为
\[
\sigma_2(T\times P):=\bigl(\cdim(T),\,\pcdim(P)\bigr).
\]
则
\[
\sigma_2\bigl(\mathcal E_6^{\mathrm{mix}}\bigr)=(3,2).
\]
并且此签名是极小的：若
\[
f:\mathcal E_6^{\mathrm{mix}}
\longrightarrow
\TT^{\,r}\times \ZZ_2^{\,s}
\]
为连续群同态，且其在 \(T_6\) 与 \(P_2\) 上的限制都忠实，则必有
\[
r\ge 3,
\qquad
s\ge 2.
\]
\end{proposition}
\begin{proof}
由定理 \ref{thm:conclusion-window6-minimal-compact-lie-completion-circle-rigidity}，\(\cdim(T_6)=3\)。另一方面，定理 \ref{thm:conclusion-window6-8adic-readout-shift-conjugacy} 与推论 \ref{cor:xi-time-part9gc-resolution-tower-rank-two-rigidity} 已把 dyadic 地址空间固定为
\[
P_2=T_2(E_{\min})\cong \ZZ_2^2,
\]
故 \(\pcdim(P_2)=2\)。于是
\[
\sigma_2\bigl(\mathcal E_6^{\mathrm{mix}}\bigr)=(3,2).
\]

若 \(f|_{T_6}\) 忠实，则 \(f(T_6)\) 是 \(\TT^{\,r}\) 的一个三维闭连通子环面，从而 \(r\ge 3\)。若 \(f|_{P_2}\) 忠实，则因紧致到 Hausdorff 的连续单射自动为到其像的同胚，\(f(P_2)\) 与 \(\ZZ_2^2\) 同构，因而
\[
2=\pcdim\bigl(f(P_2)\bigr)\le \pcdim\bigl(\ZZ_2^{\,s}\bigr)=s.
\]
故 \(s\ge 2\)。这正说明 \((3,2)\) 不能在保持两侧忠实性的前提下再被压缩。证毕。
\end{proof}

\begin{corollary}[物理骨架的强制三分层]\label{cor:conclusion-window6-physical-skeleton-forced-stratification}
若把 window-\(6\) 的可见结构解释为一个边界--体相耦合系统，则与上述定理链相容的极小运动学骨架被强制为
\[
T_6\curvearrowright S_{\mathrm{spin}}\times P_2,
\]
其中 \(T_6\cong \TT^3\) 承载连续相位，\(S_{\mathrm{spin}}\) 为定理 \ref{thm:conclusion-window6-minimal-fiber-spin-hypercube} 中承载 \(Y_2\) 的八维自旋模，\(P_2\cong \ZZ_2^2\) 承载 dyadic 分辨率地址。换言之，window-\(6\) 的最小统一对象既不是单一环面，也不是单一有限超立方体，而是一个“三环面相位--八重自旋边界层--秩二 dyadic 地址账本”的耦合骨架。
\end{corollary}
\begin{proof}
定理 \ref{thm:conclusion-window6-minuscule-eightfold-b3-spin7} 把最小 Lie 型唯一压缩到 \(\mathrm{Spin}(7)\)，定理 \ref{thm:conclusion-window6-minimal-fiber-spin-hypercube} 把最小纤维层提升为自旋权超立方体，而命题 \ref{prop:conclusion-window6-dyadic-mixed-signature-minimality} 说明 dyadic 地址轴不可再降地保留一个秩二 pro-\(2\) 账本。将这三部分合并，即得所述三分层骨架。证毕。
\end{proof}

\begin{conjecture}[minuscule 边界刚性]\label{conj:conclusion-minuscule-boundary-rigidity}
对任意分辨率窗口 \(m\)，设存在刚性分裂
\[
X_m=X_m^{\mathrm{cyc}}\sqcup X_m^{\mathrm{bdry}},
\]
并记
\[
Y_m^{\min}:=\{x\in X_m:\ d_m^{\mathrm{bin}}(x)=d_{\min}\}.
\]
若满足：
\begin{enumerate}
  \item 循环层 \(X_m^{\mathrm{cyc}}\) 允许一个不可约的极小紧 Lie 完成；
  \item 最小纤维层 \(Y_m^{\min}\) 可在该完成的某个非零 minuscule 表示中实现为单一 Weyl 权轨道；
\end{enumerate}
则极小 Lie 型应由二元计数
\[
\bigl(|X_m|,\ |Y_m^{\min}|\bigr)
\]
唯一决定（以 Langlands 对偶为模）；并且由 \(X_m^{\mathrm{bdry}}\) 所生成的模二 anomaly 子格应当规范同构于该完成的极大环面二挠
\[
T_m[2].
\]
对 window-\(6\)，该猜想的首个非平凡实例即为
\[
\bigl(|X_6|,\ |Y_2|\bigr)=(21,8)
\qquad\Longrightarrow\qquad
\mathrm{Spin}(7)
\]
（在单连通侧）。
\end{conjecture}

\begin{conjecture}[boundary 八重超选作为 \texorpdfstring{\(\mathrm{Spin}(7)\)}{Spin(7)} 自旋权的模二阴影]\label{conj:conclusion-window6-boundary-superselection-spin7-mod2-shadow}
记 \(\Sigma_8\) 为定理 \ref{thm:conclusion-window6-minuscule-eightfold-b3-spin7} 中 \(\mathrm{Spin}(7)\) 自旋表示的八个权。则应存在一个与 boundary 三轴循环相容的规范识别
\[
\widehat Z_6^{\mathrm{bd}}
\cong
\Sigma_8 \pmod 2,
\]
使得定理 \ref{thm:conclusion-window6-boundary-superselection-c3-orbit-stratification} 中的刚性分裂
\[
(1+1+3+3)
\]
正是 \(\Sigma_8\) 在该循环作用下的精确轨道分解。换言之，window-\(6\) 的八重超选扇区不应被看作偶然的八点角色立方体，而应理解为唯一 \(B_3\)-minuscule 完成在模 \(2\) 阴影中的严格残迹。
\end{conjecture}
\begin{proof}
定理 \ref{thm:conclusion-window6-minuscule-eightfold-b3-spin7} 已把最小 Lie 型唯一压缩到 \(\mathrm{Spin}(7)\)，并把最小纤维层 \(Y_2\) 识别为自旋表示的八个权。另一方面，定理 \ref{thm:conclusion-window6-boundary-torus-2torsion} 与定理 \ref{thm:conclusion-window6-pinned-datum-toral-completion-eightsector} 已把 boundary parity 识别为极大环面的全部二挠，并给出严格的八重超选分解；再由定理 \ref{thm:conclusion-window6-boundary-superselection-c3-orbit-stratification}，这八个角色在 boundary 三轴循环下并非均匀分布，而是刚性地裂解为
\[
1+1+3+3
\]
四个轨道。

因此，最自然的统一解释是：\(\widehat Z_6^{\mathrm{bd}}\) 应当正是 \(\Sigma_8\) 的模二阴影，而上述四轨道分裂对应于自旋权系在同一三轴循环作用下的精确轨道结构。当前尚缺的是把这一识别写成与 pinned datum 兼容、且不依赖额外坐标选择的规范构造，故保留为猜想。证毕。
\end{proof}

\paragraph{统一读法}
上述链条说明：伴随层的 \(21=18+3\) root--Cartan 完成把 window-\(6\) 的连通可见圆维锁定为 \(3\)；辅助 minuscule 层的八个最小纤维则把 Langlands 对偶的二义性压缩为 \(\mathrm{Spin}(7)\)，并将最小层提升为自旋权超立方体；地址层的 \(T_2(E_{\min})\cong \ZZ_2^2\) 则给出不可再降的 dyadic 账本。于是几何上出现的是一个八点自旋超立方体，代数上出现的是 rank-\(3\) 的极小紧 Lie 完成，矩阵上同时进入 \(\mathfrak{so}(7)\) 的 \(7\times 7\) 斜对称模型与其 \(8\) 维自旋模，而信息上则形成不可互换的 \((3,2)\) 混合签名。相应的最小物理读法，就是连续相位、边界自旋层与 dyadic 地址账本三者的耦合。

\endinput
